%!TEX root = main.tex
%!TEX spellcheck = en_US

\section{Introduction}
\label{sec:intro}

Ensuring the stability, security, and performance of network infrastructures
hinges on careful management when introducing a new routing device or
modifying the configuration of an existing device. Even a single device's
misconfiguration can lead to black holes, unintended network access,
inefficient routes, and a range of other performance and security
vulnerabilities. For instance, a misconfigured access control list (ACL) may
unintentionally allow unauthorized traffic, exposing the device (and potentially
the broader network) to malicious attacks; a buggy routing policy may create
loops, thereby degrading performance and rendering local services inaccessible.

Researchers and practitioners have developed a variety of tools to detect
misconfigurations in routing devices. On one hand, \textbf{model
	checkers}~\cite{fogel2015general, beckett2017general, abhashkumar2020tiramisu,
	prabhu2020plankton, zhang2022sre, steffen2020netdice, ye2020hoyan,
	ritchey2000using,al2011configchecker, jeffrey2009model} rely on manually
defined \emph{semantics}, which we define here as \emph{the set of correctness
	constraints, protocol-specific rules, and policy-specific requirements that
	configurations must satisfy}. These semantics are encoded in policies or rules
checked against the router's configuration. The advantage of model checkers
lies in their ability to encode deeper, context-specific knowledge: indeed,
for different configuration types, software versions, or vendor-specific
features, one can tailor the rules to accurately capture correctness. The key
limitation, however, is the lack of automation. The reliability of these
policy-based tools depends heavily on the correctness and completeness of the
manually inferred semantics. Moreover, for evolving device types, a variety of
configuration versions, or large sets of heterogeneous setups, the labor
required to add, remove, or modify entire sets of policy rules quickly becomes
prohibitive.

On the other hand, \textbf{consistency
	checkers}~\cite{kakarla2024diffy,kakarla2020finding, le2006minerals,
	feamster2005detecting, tang2021campion, le2008detecting, le2006characterization}
compare one device’s configuration with those of peer devices, different
snapshots of the same device, or best-practice baselines to flag
inconsistencies. This approach lends itself to automation because statistical
or pattern-based checks require minimal human intervention in crafting
explicit policies. However, consistency alone does not necessarily imply
correctness. A device’s configuration may align perfectly with peer devices or
earlier snapshots yet still harbor protocol-level or internal conflicts.

\smallskip
\noindent
\textbf{Research Question:}
\emph{Is it possible to automate the extraction and analysis of \emph{semantics}
	within a routing device’s configuration such that we combine the rich
	semantic understanding of model checkers (given sufficient and relevant
	context) with the automated flexibility of consistency checkers? In doing so,
	can we minimize the burdensome process of manually crafting policies?}

\smallskip

As a promising step toward addressing this question, researchers are beginning
to use large language models (LLMs) to parse and analyze configuration files.
Recent \textbf{LLM-based Q\&A
	tools}~\cite{bogdanov2024leveraging,chen2024automatic,wang2024identifying,liu2024large,
	wang2024netconfeval, lian2023configuration} feed entire or partial
configuration lines into pre-trained transformer models, allowing these LLMs
to detect syntax issues and some category of semantic errors. When provided
with appropriate sections of configuration (\eg, parameters relevant to an ACL
block), an LLM can often derive correct semantic constraints from its
pre-trained corpus and identify subtle errors.

All of these techniques—model checkers, consistency checkers, and LLM-based
Q\&A tools—rely on the notion of \emph{context}. Model checkers require
knowledge of protocols, device configurations, and forwarding policies;
consistency checkers require comparative baselines; and LLMs need sufficiently
rich prompts containing relevant configuration lines. Crucially, the
\emph{correctness} and \emph{completeness} of the context provided directly
affect the quality of the detection. For instance, overlooking key ACL lines
may conceal particular misconfigurations, while feeding in irrelevant or
incorrect lines may spawn false alarms.

Assembling high-quality context, however, can be a substantial burden.
Providing detailed protocol semantics, capturing vendor- or version-specific
syntax, and enumerating idiosyncratic device features is both time-consuming
and error-prone, even in a single-router scenario. Fortunately, LLMs hold
significant promise here because they are typically pre-trained on vast
corpora of network-related texts, including protocol specifications, vendor
documentation, and publicly available configurations. This \emph{generic
	configuration context} can be supplemented by \emph{device-specific context}
(\ie, the precise lines of the configuration under analysis). In principle, if
the LLM is given enough relevant device-specific data, it may automatically
\emph{derive} many of the semantics needed to check for misconfigurations.

Yet simply presenting an entire configuration in one massive prompt can lead
to \emph{context overload}, which may dilute the model's focus and exceed token
limits. Conversely, splitting a configuration into multiple prompts can cause
the LLM to lose track of dependencies across different parts of the
configuration, hindering its ability to detect correlated errors.

In this paper, we address the problem of precisely extracting and integrating
network-specific context for a given routing device into LLM prompts to enhance
the detection of router misconfigurations. Specifically, we introduce the
\emph{Context-Aware Iterative Prompting} (\sysname{}) framework, which tackles
three key challenges:

\mypara{Challenge 1: Automatic and Accurate Context Mining} Router
configurations tend to be long and
complex~\cite{benson2009complexitymetrics}, featuring lines that mutually
interact in nontrivial ways. The task is to automatically extract and select
only the most relevant context to limit computational expense and reduce the
introduction of irrelevant information that might degrade detection accuracy.
\textbf{Solution:} We model the configuration as a hierarchical tree, where
each line corresponds to a path from the root to a parameter value. This
structure enables systematic identification of three main categories of
context---\emph{neighboring}, \emph{similar}, and \emph{referenced}
configuration lines—ensuring the LLM receives a concise yet sufficient corpus
for accurate analysis.

\mypara{Challenge 2: Parameter-Value Ambiguity in Referenceable Context}
Configurations contain both pre-defined and user-defined parameters, and the
latter often require specialized context to interpret correctly. For example,
an interface might reference a user-defined ACL whose definition resides
elsewhere in the configuration. Incorrectly labeling or omitting these
references can result in extraneous context and inflate computational costs.
\textbf{Solution:} We employ an existence-based labeling approach within the
configuration tree to distinguish user-defined values by detecting their
occurrences as internal nodes in other paths. We further refine these
classifications with a majority-voting scheme, ensuring values are treated
consistently across multiple contexts.

\mypara{Challenge 3: Managing Context Overload in Prompts} Even after
precisely extracting relevant information, the quantity of required context
can exceed the capacity of a single prompt. Overloading an LLM with excessive
detail can reduce coherence and degrade performance. \textbf{Solution:} In our
iterative prompting framework, the LLM explicitly requests further context
when needed, incrementally refining the prompt and focusing on the most
critical lines. This approach mitigates the risk of overload, enabling more
precise misconfiguration detection.

We evaluate \sysname{} through two distinct case studies. First, we take
real-world device configuration snapshots, inject a series of realistic
updates (including both correct modifications and latent misconfigurations),
and show that \sysname{} outperforms existing partition-based LLM approaches,
model checkers, and consistency checkers by at least 30\% in accurate error
detection. Second, we run \sysname{} exhaustively on configuration snapshots
from a medium-sized campus network. \sysname{} identifies more than 20
previously overlooked misconfigurations, further underscoring its practical
utility.
